\documentclass[11pt]{article}
\usepackage[a4paper,margin=23mm]{geometry}
\usepackage{amsmath,amssymb,amsthm,booktabs,longtable,hyperref}
\usepackage{microtype}\emergencystretch=2em
\newtheorem{theorem}{Theorem}\newtheorem{lemma}[theorem]{Lemma}
\newcommand{\GL}{\operatorname{GL}}\newcommand{\SL}{\operatorname{SL}}
\newcommand{\PGL}{\operatorname{PGL}}\newcommand{\PSL}{\operatorname{PSL}}
\newcommand{\diag}{\operatorname{diag}}\newcommand{\F}{\mathbb F}
\title{Three involutions in finite general linear groups}
\author{A complete principal-result recovery of Markovskaya and Nuzhin (2025)}
\date{}
\begin{document}\maketitle
We recover the two classifications of I.~A.~Markovskaya and Ya.~N.~Nuzhin
[MN], retaining their explicitly borrowed classifications and integral
generation theorem. Their constructive matrices are retained and their
root-generation steps are supplied by the complete common paper
\emph{Root parameters and fixed-space obstructions} [U]. We correct three
displayed matrix expressions below. These corrections do not contradict
either principal classification. The intended domain is $n\geq2$, as in
both source tables; the literal omission of that domain in the theorem
sentences is recorded separately.

\section{Definitions and retained inputs}
Call a group admissible if it is generated by three elements of order at
most two, the first two commuting. For a nontrivial group this is equivalent
to the source convention allowing repeated involutions: an identity entry
can be replaced by an existing suitable involution (if a commuting entry
is identity, repeat the other commuting entry; if both are identities,
repeat the remaining entry). The trivial group is included by convention.
Thus all dihedral groups are admissible. Quotients are admissible: simply
take the images of the marked generators. For $H\lhd G$, a tuple generates
$G/H$ exactly when it generates $G$ together with $H$, by taking the full
inverse image of the generated quotient subgroup.

The following inputs are retained, not reclassified in this paper.
\begin{enumerate}
\item The known classification of admissible finite simple linear groups,
as used in [MN, introduction and proof of Theorem 2]: among $\PSL_n(q)$,
$n\geq2$, the exclusions are $\PSL_2(7)$, $\PSL_2(9)$, every
$\PSL_3(q)$, and every $\PSL_4(2^a)$. The small nonsimple groups
$\PSL_2(2)$ and $\PSL_2(3)$ are handled separately if needed. We use this
classification only for $n\geq3$; it is attributed to Nuzhin [N19].
\item $\PGL_2(q)$ is admissible for every prime power $q>2$, by
Sjerve--Cherkassoff [SC]. The remaining $\PGL_2(2)\cong S_3$ is dihedral.
\item $\GL_n(\mathbb Z)$ is admissible for $n\geq5$, by the earlier
Markovskaya--Nuzhin theorem [MN23]. Its proof is retained, not replaced by
computing finitely many reductions.
\item The rank-three intersection criterion of Conder--Oliveros [CO,
Corollary 4.2]: three distinct involutions, two commuting, in a finite
group with no nontrivial cyclic normal subgroup give a regular abstract
rank-three polytope. The standard group--polytope correspondence is
retained with this criterion.
\end{enumerate}
These are genuine dependencies, whose complete external proof and
software costs are unknown here. The source's general rank-two
transvection statements [MN, Lemma 10], attributed to Levchuk [L83b], are
also retained as background statements, but none is a premise of our
constructive proof. All root and fixed-space steps actually used below
are proved in [U].
For precision the three retained rank-two assertions are: (a) if $q\ne9$,
$V\subseteq\F_q^*$ has at least two elements and contains an element
lying in no proper subfield, then $\langle t_{12}(V),t_{21}(1)\rangle
=\SL_2(q)$; (b) for nonzero $a,b$, the group
$\langle t_{12}(a),t_{21}(b)\rangle$ contains some
$\left(\begin{smallmatrix}0&-w^{-1}\\w&0\end{smallmatrix}\right)$ and
its square $-I$; (c) for odd $q$ and primitive $u$,
$\langle t_{12}(u),t_{21}(1)\rangle=\SL_2(q)$.
These retain the source's exact hypotheses and are not claimed as
consequences of our recognition theorem, which requires at least three
coordinates.

\section{Determinants, quotients, and the linear classification}
For a commutative unital ring $K$, the determinant has the section
$a\mapsto\diag(a,1,\ldots,1)$, and hence
\[\GL_n(K)=\SL_n(K)\rtimes\{\diag(a,1,\ldots,1):a\in K^*\}.\]
This argument does not require elementary generation over $K$.
If $K^*$ is finite and $\gcd(n,|K^*|)=1$, its $n$th-power map is an
automorphism. Every determinant is therefore the determinant of a scalar
matrix, with trivial intersection between scalars and $\SL_n(K)$. Thus
$\GL_n(K)=\SL_n(K)\times K^*I$ and
$\PGL_n(K)\cong\SL_n(K)=\PSL_n(K)$.

Over a finite field, Gaussian elimination shows that all elementary
transvections generate $\SL_n(q)$. Reduction from $\GL_n(\mathbb Z)$
modulo a prime $p$ contains all these transvections and
$\diag(-1,1,\ldots,1)$; its image is exactly the determinant-$\{1,-1\}$
subgroup of $\GL_n(p)$. It is onto precisely for $p=2,3$. For $n=1$
the same determinant description also holds, though that dimension is
not in the present classifications.

\begin{theorem}\label{gl}
For $n\geq2$, the group $\GL_n(q)$ is admissible exactly when
\[
 (q=2\text{ and }(n=2\text{ or }n\geq5))
 \quad\text{or}\quad(q=3\text{ and }n\geq5).
\]
\end{theorem}
\begin{proof}
For $q\geq4$, determinants of involutions lie in $\{1,-1\}$, a proper
subgroup of $\F_q^*$, so even unrestricted involution generation fails.
For $q=2$, all four linear/projective/special groups coincide. The retained
simple classification excludes $n=3,4$ and allows $n\geq5$; $n=2$ gives
$S_3$. For $q=3,n\geq5$, reduce the retained integral generating triple.
For $n=4,q=3$, [U, symmetric-square obstruction] supplies the full
ten-dimensional argument. For $n=3,q=3$, the direct product just proved
has $\SL_3(3)=\PSL_3(3)$ as a quotient, which is excluded by the retained
classification.

It remains to exclude $\GL_2(3)$. Commuting involutions are simultaneously
diagonalizable, so place the commuting pair in the diagonal group.
Every nonscalar involution is
$\left(\begin{smallmatrix}a&b\\c&-a\end{smallmatrix}\right)$ with
$a^2+bc=1$. If $a=\pm1$, then $bc=0$, so it is upper or lower triangular;
if $a=0$, then $bc=1$, so it is monomial. A diagonal group together with
one such matrix remains in the corresponding triangular or monomial
proper subgroup. Scalar possibilities cannot enlarge this subgroup.
\end{proof}

The determinant quotient of $\PGL_n(q)$ is
\[\PGL_n(q)/\PSL_n(q)\cong\F_q^*/(\F_q^*)^n\cong C_{\gcd(n,q-1)}.
\]
An abelian group generated by involutions has exponent at most two. Hence
if $\gcd(n,q-1)>2$, $\PGL_n(q)$ is not generated by any collection of
involutions. This proves [MN, Lemma 11], including its stated divisible
special case, without restrictions on lifts of projective involutions.

\section{Common coordinates for even dimensions}
For the remaining case put $n=2k\geq4$, $q$ odd, and
$\gcd(n,q-1)=2$. Choose a primitive $u\in\F_q^*$ and set
$\epsilon=(-1)^k$. All matrices act on column vectors. Let
\begin{gather*}
 \mu e_i=e_{i+1}\quad(i\bmod n),\qquad \tau e_i=e_{n+1-i},\\
 B=\diag(\underbrace{\epsilon/u,\ldots,\epsilon/u}_{k},
                    \underbrace{1,\ldots,1}_{k})\tau,
 \qquad C=\tau\mu,\quad E=BC.
\end{gather*}
Then $B^2=(\epsilon/u)I$, $C^2=I$, and $E$ is the displayed diagonal
matrix times $\mu$. With cyclic indices and its diagonal entries $d_i$,
\begin{equation}\label{transport}
 t_{ij}(a)^E=t_{i+1,j+1}(a\,d_{i+1}/d_{j+1}).
\end{equation}
The conjugation convention throughout is $X^Y=YXY^{-1}$.
Directly,
\[\det B=u^{-k},\qquad\det C=(-1)^{k-1}.\]
If $k$ is odd, $\gcd(k,q-1)=1$, so $\det B$ is primitive. If $k$ is
even, write $q-1=2m$ and $k=2h$. The gcd assumption gives $m$ odd and
$\gcd(h,m)=1$; $\det(BC)=-u^{-k}=u^{m-2h}$ is primitive, since
$\gcd(m-2h,2m)=1$. Thus once $\SL_n(q)$ lies in $\langle A,B,C\rangle$,
this group is the whole $\GL_n(q)$ in the cases using these coordinates.

\subsection{All dimensions at least eight}
Let $A$ have diagonal entries $-1$ at $2,3,n-2,n-1$, all other diagonal
entries $1$, and just two off-diagonal entries
$A_{21}=A_{n-1,n}=1$. Then $A^2=I$ and $AB=BA$.
Define the following words without changing their order:
\[
 X=[A,A^E],\quad T=((XA^{E^2})^2)^E,\quad Z=[T,X],\quad
 R=[A,Z],\quad V=R^B,\quad W=[Z,V].
\]
The complete nonidentity entries of these words are as follows; products
in each table entry have commuting factors.
\begin{center}\small
\begin{tabular}{c|l}
word&value\\\hline
$X$&$t_{31}(-1)t_{n-1,1}(\epsilon u)$\\
$T$, $n=8$&$t_{43}(-u)t_{52}(u)t_{53}(-u^2)t_{83}(u^3)$\\
$T$, $n\geq10$&$t_{43}(-\epsilon u)t_{52}(1)t_{53}(-\epsilon u)t_{n3}(\epsilon u^3)$\\
$Z$, $n=8$&$t_{41}(u)t_{51}(u^2)t_{81}(-u^3)$\\
$Z$, $n\geq10$&$t_{41}(\epsilon u)t_{51}(\epsilon u)t_{n1}(-\epsilon u^3)$\\
$R$&$t_{n-1,1}(-\epsilon u^3)$\\
$V$&$t_{2n}(-u^2)$\\
$W$&$t_{21}(-\epsilon u^5)$
\end{tabular}
\end{center}
These identities hold as Laurent-polynomial identities over
$\mathbb Z[u,u^{-1}]$, with $\epsilon=(-1)^k$. One way to verify them
uniformly is to conjugate the diagonal entries and the two off-diagonal
entries of $A$ using \eqref{transport}, and multiply using
$E_{ij}E_{ab}=\delta_{ja}E_{ib}$. Only indices
$1,2,3,4,5,n-2,n-1,n$ occur. The distinction in the two $T,Z$ rows is
exactly that $d_5=1$ for $n=8$, whereas $d_5=\epsilon/u$ for $n\geq10$.
No enumeration of dimensions or fields is a premise of these identities.

The nonzero adjacent root $W$ and the monomial cycle $E$ reach every
root by [U, root propagation]. Apply [U, finite root-parameter recognition]
to $M=\langle A,B,C\rangle$. Its root-generated subgroup contains the
determinant-one diagonal sign matrices. Isolating the second factor of
$X$ by [U, isolation] gives $t_{n-1,1}(\epsilon u)$. Together with $R$
this puts $-u^2$ in the common parameter subfield; thus that subfield is
the whole field, by [U, primitive-square test]. Hence $M$ contains
$\SL_n(q)$ and is $\GL_n(q)$ by the determinant argument. The projective
images of $A,B,C$ are involutions, the first two commute, and they
generate $\PGL_n(q)$.

\subsection{Dimension six}
For $q=3$ Theorem~\ref{gl} already supplies the required projective triple.
Otherwise choose the primitive element so that $h=u^2+4u\ne0$.
This is possible: the only forbidden nonzero value is $-4$, and a cyclic
group of even order $q-1\geq4$ has at least two primitive elements.
Keep $B,C,E$ above, and let $A$ have diagonal entries $-1$ at $1,6$,
$1$ elsewhere, and $A_{21}=A_{56}=-1$. Again $A^2=I$ and $AB=BA$.
The exact words are
\begin{align*}
 X=[A,A^E]&=t_{31}(-1)t_{51}(-u),\\
 Q=[A^{E^4},X]&=t_{51}(2u)t_{61}(-u^2),\\
 [A,Q]&=t_{51}(-h),\\
 (A^{E^4}t_{51}(h))^2&=t_{61}(uh),\qquad
 t_{61}(uh)^E=t_{12}(uh).
\end{align*}
These follow from \eqref{transport} and the same matrix-unit product rule.
Since $uh\ne0$, propagation again gives all roots. Root recognition then
supplies $D=\diag(1,1,1,1,-1,-1)$, and
\[
 [D,X]=t_{51}(2u),\qquad t_{51}(2u)^C=t_{15}(2u).
\]
Their opposite-root parameter product is $4u^2$; the primitive-square
test gives the full field. Thus $M$ contains $\SL_6(q)$, and
$\det B=u^{-3}$ is primitive because $\gcd(6,q-1)=2$.
This proves the six-dimensional branch without assuming that $2u$ itself
is a primitive multiplicative element.

\section{Dimension four}
Here $\gcd(4,q-1)=2$ means $q\equiv3\pmod4$, so $-1$ is nonsquare.
In both branches put
\[
 A=\begin{pmatrix}0&0&0&-1\\0&0&-1&0\\0&1&0&0\\1&0&0&0\end{pmatrix}.
\]
We will obtain $\SL_4(q)\leq M=\langle A,B,G\rangle$, with
$\det A=\det B=1$ and $\det G=-1$. It follows that $M$ is exactly the
determinant-$\{1,-1\}$ subgroup. Its projective image is all of
$\PGL_4(q)$ because $\F_q^*/(\F_q^*)^4$ has order two and $-1$ represents
its nontrivial coset.

\subsection{Characteristic three}
Take
\[
 B=\begin{pmatrix}0&0&-1&0\\0&0&0&1\\1&0&0&0\\0&-1&0&0\end{pmatrix},
 \qquad
 G=\begin{pmatrix}-1&0&0&0\\-1&1&0&0\\-u&0&1&0\\-u&0&0&1\end{pmatrix}.
\]
Then $A^2=B^2=-I$, $G^2=I$, and $AB=-BA$. Thus their projective images
are involutions, the first two commuting. Set
\[
 J=GG^{AB}G,\qquad D=(G^{AB}G)^2J^B,\qquad
 T=GDD^AG^A,\qquad P=(T^BJ)^2.
\]
Multiplication over $\F_3[u]$ gives
\begin{align*}
 D&=\begin{pmatrix}-1&0&0&0\\0&-1&u&u\\0&u&-u^2-1&1-u^2\\-u&0&0&1\end{pmatrix},\\
 DD^A&=\begin{pmatrix}-1&0&0&-u\\-1&1&0&-u\\-u&0&1&u^2-1\\-u&0&0&-u^2-1\end{pmatrix},\\
 T&=t_{24}(u)t_{34}(u^2),\qquad P=t_{42}(w),\qquad w=u^3+u,\\
 P^{AB}&=t_{31}(w),\quad P^A=t_{13}(-w),\quad
 [G^{AB},P^{AB}]=t_{32}(w),\quad t_{32}(w)^A=t_{23}(-w).
\end{align*}
Here $w\ne0$, because $u^2=-1$ is impossible. The four displayed roots
on coordinates $1,2,3$ give a strongly connected directed graph. Applied
to the subgroup they generate on these three coordinates, root recognition
supplies determinant-one sign matrices and monomial lifts of $S_3$ on
these coordinates, fixing the fourth coordinate. In particular
\[
 (T\diag(-1,1,-1,1))^2=t_{24}(-u),\qquad
 t_{24}(-u)^B=t_{42}(u).
\]
The permutation parts of that $S_3$ and of $AB$, namely $(12)(34)$,
generate $S_4$. Hence conjugating a nonzero root reaches all ordered
coordinate pairs. Root recognition in dimension four and the opposite
parameters $u,u$ now give $\SL_4(q)\leq M$.

The additional source assertion about linear lifts is also valid:
$\langle uA,B,G\rangle=\GL_4(q)$. Write $q-1=2m$, $m$ odd.
Since $(uA)^2=-u^2I=u^{m+2}I$ and
$\gcd(m+2,2m)=1$, this group contains every scalar, hence $A$ and $M$.
Its determinants generate $\langle-1,u^4\rangle=\F_q^*$.

\subsection{Odd characteristic different from three}
Take instead
\[
 B=\begin{pmatrix}0&0&1&0\\0&0&0&1\\1&0&0&0\\0&1&0&0\end{pmatrix},
 \qquad
 G=\begin{pmatrix}-1&0&0&0\\1&1&0&0\\u&0&1&0\\u&0&0&1\end{pmatrix}.
\]
Here $A^2=-I$, $B^2=G^2=I$, and $AB=-BA$. Let $J=GG^{AB}G$.
Its upper-left block is
$\left(\begin{smallmatrix}0&-1\\-1&0\end{smallmatrix}\right)$ and its
lower block is $I_2$. The products
$JAB$ and $(JAB)^A$ are the ordinary swaps $(34)$ and $(12)$ respectively.
Consequently $D=(JAB)^AJ=\diag(-1,-1,1,1)$ lies in $M$, and
\[
 [D,G]=t_{31}(2u)t_{41}(2u).
\]
Taking an integer power with exponent representing $-1/2\in\F_p$ gives
$V=t_{31}(-u)t_{41}(-u)$. Direct multiplication gives
\[
 E_0=GV=\diag(-1,1,1,1)t_{21}(1),\qquad
 U=V^A=t_{14}(u)t_{24}(u),\qquad
 (UE_0)^2=t_{24}(3u).
\]
As $3\ne0$, an integer power gives $t_{24}(u)$ and conjugation by $B$
gives $t_{42}(u)$. The monomial matrices $A,B,(12),(34)$ preserve the
partition $\{1,2\}\mid\{3,4\}$ and transport this root onto every
ordered pair across the partition. This directed complete bipartite
graph is strongly connected. Propagation gives all roots, and the
opposite parameters give the full field, proving $\SL_4(q)\leq M$.
This supplies the source's omitted last generation step without
incorrectly assuming that these monomial permutation parts form $S_4$.

\section{Projective classification and both tables}
\begin{theorem}\label{pgl}
For $n\geq2$, $\PGL_n(q)$ is admissible precisely in these cases:
\[
 n=2;\qquad n\geq4,\ \gcd(n,q-1)=2;\qquad
 n\geq5,\ \gcd(n,q-1)=1.
\]
\end{theorem}
\begin{proof}
The degree-two case is retained from [SC], with $q=2$ handled directly.
The determinant quotient excludes gcd greater than two. For gcd one,
$\PGL_n(q)\cong\PSL_n(q)$, so the retained simple classification gives
exactly the asserted possibilities: in degree four gcd one forces $q$
even, and those groups are excluded. For degree three the gcd is either
one or three, giving exclusion in both cases. Gcd two forces even
$n\geq4$ and odd $q$; the preceding three dimension branches of the
construction (four, six, and at least eight) exhaust this case.
\end{proof}

\begin{center}
\begin{tabular}{lll}\toprule
group&parameters&admissible\\\midrule
$\GL_n(q)$&$n=2,q=2$&yes\\
&$n=2,q\ne2$&no\\
&$n=3$, any $q$&no\\
&$n=4$, any $q$&no\\
&$n\geq5,q=2,3$&yes\\
&$n\geq5,q\ne2,3$&no\\\midrule
$\PGL_n(q)$&$n=2$, any $q$&yes\\
&$n=3$, any $q$&no\\
&$n\geq4,\gcd(n,q-1)>2$&no\\
&$n\geq4,\gcd(n,q-1)=2$&yes\\
&$n=4,\gcd(n,q-1)=1$&no\\
&$n\geq5,\gcd(n,q-1)=1$&yes\\\bottomrule
\end{tabular}
\end{center}
Thus every row of the two source tables follows, with the explicitly
retained inputs specified above.

Every positive projective group other than $\PGL_2(2)$ also gives the
claimed rank-three regular polytope. The generators can be taken distinct:
a repeated pair would make the group dihedral, whereas these groups are
not dihedral. For the almost simple cases, a cyclic normal subgroup $N$
would intersect the simple socle trivially and therefore centralize it,
since $[N,S]\leq N\cap S$. The centralizer of $\PSL_n(q)$ in
$\PGL_n(q)$ is trivial: a projective centralizer fixes each elementary
transvection up to a scalar, that scalar must be one by its sole eigenvalue
$1$, and commuting with all matrix units forces a scalar matrix.
The exceptional small positive group $\PGL_2(3)\cong S_4$ has no
nontrivial cyclic normal subgroup either, as its normal subgroups are
$1,V_4,A_4,S_4$. Apply the retained [CO] criterion. No new answer to the
broad problems mentioned in the source's survey citation is claimed.

\section{Literal-source qualifications and attribution}
\begin{enumerate}
\item The source's theorem sentences omit $n\geq2$, while the two tables
begin at degree two. Without that intended domain, the sentences need
correction: $\PGL_1(q)$ is trivial for every $q$, and $\GL_1(2)$ and
$\GL_1(3)$ are admissible. No excluded degree-$\geq2$ case is changed.
\item On printed p.~254 the intermediate determinant expression for $B$
has a sign inconsistency when $k$ is even. The correct value is
$u^{-k}$, as computed from the full matrix; the primitive determinant
argument above covers both parities.
\item On p.~256 the $(4,4)$ entry of $DD^A$ is $-u^2-1$, not $-u-1$.
The subsequent source word for $T$ agrees with the corrected entry.
\item On p.~257, in characteristic different from three, the conjugator
in $t_{14}(u)t_{24}(u)=V^A$ must be $A$, not the printed $G$:
$V^G=t_{31}(u)t_{41}(u)$. The next root identity requires the square of
the whole product $(UE_0)^2$; squaring only $E_0$ leaves $U$ unchanged.
These are explicit word corrections, not counterexamples to Theorem 2.
\item The source's root-recognition applications are completed here by
forming the root-generated subgroup and isolating paired root factors.
In dimension six the opposite-product $4u^2$ argument avoids assuming
that $2u$ is multiplicatively primitive.
\item The English end matter prints received/revised dates in 2024,
whereas the Russian end matter gives 2025. This is bibliographic only.
\end{enumerate}
All constructive matrices and the symmetric-square obstruction are due
to the source authors; our contribution is the complete elementary common
proof, explicit parameter extraction and corrected identities. The older
private rank-three diagonal-extension work uses different marked groups
and is not silently imported here. Symbolic checks of displayed words
are optional diagnostics: the written matrix identities, multiplication
rule and uniform arguments are the mathematical proof. No new Lean
declaration, finite search premise, compression ratio, or conjecture
resolution is asserted.

\begin{thebibliography}{9}\small
\bibitem{MN}[MN] I.~A.~Markovskaya and Ya.~N.~Nuzhin,
On generation of the groups $\GL_n(q)$ and $\PGL_n(q)$ by three involutions,
two of which commute, Trudy Inst. Mat. Mekh. UrO RAN 31(4) (2025),
247--259. DOI 10.21538/0134-4889-2025-31-4-fon-03.
\bibitem{U}[U] Root parameters and fixed-space obstructions, complete
companion manuscript to this recovery.
\bibitem{N19}[N19] Ya.~N.~Nuzhin, Generating sets of involutions of finite
simple groups, Algebra and Logic 58 (2019), 288--293.
\bibitem{SC}[SC] D.~Sjerve and M.~Cherkassoff, On groups generated by three
involutions, two of which commute, CRM Proc. Lecture Notes 6 (1994),
169--185.
\bibitem{MN23}[MN23] I.~A.~Markovskaya and Ya.~N.~Nuzhin, On generation
of the groups $\GL_n(\mathbb Z)$ and $\PGL_n(\mathbb Z)$ by three
involutions, two of which commute, J. Siberian Federal Univ.
Mathematics \& Physics 16(4) (2023), 413--419.
\bibitem{CO}[CO] M.~Conder and D.~Oliveros, The intersection condition
for regular polytopes, J. Combin. Theory Ser. A 120 (2013), 1291--1304.
\bibitem{L}[L83b] V.~M.~Levchuk, Remark on a theorem of L.~Dickson,
Algebra and Logic 22 (1983), 306--316.
\end{thebibliography}
\end{document}
